% ==============================================================================
% capitulos/01-introduccion.tex - Introducción del Proyecto
% ==============================================================================

\chapter{Introducción}
\begin{flushleft}
\justifying
El Instituto Politécnico Nacional (IPN), institución rectora de la educación tecnológica en México, establece en su Ley Orgánica \cite{1} y su Reglamento Interno \cite{2} el compromiso de garantizar un entorno de respeto, equidad y legalidad para su comunidad. Ante la identificación de vulneraciones a los derechos de estudiantes, docentes y personal administrativo, el Instituto formalizó la creación de la Defensoría de los Derechos Politécnicos (DDP) mediante su acuerdo de creación oficial \cite{3} . Este órgano, dotado de independencia funcional, constituye el mecanismo institucional responsable de la gestión y seguimiento de inconformidades derivadas de actos u omisiones que contravengan la normativa politécnica.

La administración de estas quejas es un proceso crítico que exige, además de sensibilidad humana, una gestión operativa de alta precisión. Actualmente, la DDP rige su funcionamiento a través del Manual de Organización \cite{4} y el Manual de Procedimientos \cite{5}, instrumentos que definen las etapas de recepción, radicación, investigación y resolución. Sin embargo, la ejecución de estos procesos enfrenta desafíos tecnológicos sustanciales; el flujo de trabajo actual depende de mecanismos manuales y herramientas de oficina descentralizadas que imposibilitan la integración sistémica de la información. Esta carencia de infraestructura digital, sumada al incremento en el volumen de solicitudes, compromete la eficiencia del personal operativo y prolonga los tiempos de respuesta institucionales.

\section{Planteamiento del Problema}
\end{flushleft}
\justifying
La DDP es el órgano institucional encargado de salvaguardar la legalidad y el respeto a los derechos politécnicos de los estudiantes, docentes y personal administrativo del IPN. Su función es fundamental para la atención y resolución de inconformidades; no obstante, ante la creciente demanda de servicios y la complejidad de los procesos operativos actuales, se ha identificado la necesidad de fortalecer su eficiencia mediante la modernización tecnológica. Bajo el flujo de trabajo vigente, se han detectado tres puntos críticos que fundamentan el desarrollo de la solución propuesta en este proyecto.

\subsection{Contexto y situación problemática}
En primer término, se identifica una vulnerabilidad crítica en la estructura operativa de la Defensoría: la desproporción entre la demanda de atención ciudadana y la capacidad del capital humano especializado. Actualmente, la función sustantiva de protección legal y resolución de controversias recae de forma exclusiva sobre un cuerpo técnico de apenas cinco abogados, quienes tienen la responsabilidad de gestionar la totalidad de las quejas declaradas procedentes dentro de una comunidad que abarca cientos de miles de integrantes.

Bajo el modelo de gestión vigente, este personal de alta especialización jurídica se ve absorbido por una densidad administrativa abrumadora. Al carecer de una plataforma centralizada, los abogados deben dedicar una fracción significativa de su jornada a tareas procedimentales, tales como el registro manual de expedientes, el seguimiento analógico de plazos procesales y la organización descentralizada de evidencias. Esta carga de trabajo no sustantiva actúa como un factor de erosión del tiempo efectivo, limitando el análisis profundo que cada caso amerita para garantizar el debido proceso.

En consecuencia, la implementación de una solución tecnológica avanzada no debe entenderse simplemente como una herramienta de apoyo, sino como un imperativo estratégico de modernización. Al automatizar los flujos de trabajo repetitivos y digitalizar la trazabilidad de los casos, se faculta a los juristas para recuperar su capacidad de enfoque en la resolución sustantiva. El objetivo es transitar de un modelo de gestión reactivo y manual hacia un ecosistema digital eficiente, donde el talento humano sea liberado de la burocracia técnica para centrarse íntegramente en la salvaguarda de los derechos de la comunidad politécnica.

\subsection{Problema a resolver}
Derivado de lo anterior, existe un reto en la eficiencia de la recepción y clasificación de solicitudes durante la etapa de \textit{``Primer Contacto''}. Ya que, al recibir información por múltiples canales descentralizados, las solicitudes suelen presentarse de forma incompleta o no son competentes para el órgano de la DDP, resaltando la falta de un asistente tecnológico para el filtrado de competencia, obligando a la revisión manual, retrasando la atención a casos de alta prioridad.

Finalmente, la gestión de documentos realizada de forma tradicional y la falta de seguimiento impactan en la transparencia del proceso. La actual dependencia de métodos de archivo físicos dificulta la consulta rápida de antecedentes. La transición a una plataforma integral no solo brindaría certidumbre al quejoso sobre el estatus de su queja, sino que permitiría al personal de la DDP la revisión de documentos y datos importantes sobre los casos.

\section{Solución Propuesta}
Se propone el desarrollo de una plataforma web integral diseñada para centralizar, automatizar y optimizar el ciclo de vida de las quejas dentro de la DDP. La solución se fundamenta en una arquitectura basada en microservicios y contempla la implementación de un módulo de asistencia inteligente en la etapa de primer contacto, el cual utiliza técnicas de PLN para asistir al usuario en la determinación preliminar de la competencia de su queja. Se permitirá la gestión basada en roles, garantizando que cada actor interactúe con la información de acuerdo con sus facultades normativas, asegurando el seguimiento total del expediente desde su recepción hasta su resolución final.

\section{Justificación}
La implementación de esta plataforma web se justifica por la necesidad estratégica de fortalecer los procesos operativos de la DDP, transformando una gestión actual basada en métodos manuales y herramientas de oficina descentralizadas hacia un flujo digital inteligente.

\subsection{Ámbito social e institucional}
En el ámbito social e institucional, el proyecto busca fortalecer varios puntos empezando por la eliminación de barreras que impiden una mejor comunicación entre la comunidad politécnica y la DDP al proporcionar un canal de seguimiento transparente que elimine la incertidumbre del quejoso. De igual forma, al automatizar las tareas administrativas, se potencia la capacidad de respuesta del personal de la Defensoría y se reduce el costo de gestión operativa, permitiendo que el cuerpo de abogados centralice sus esfuerzos en el análisis jurídico de fondo a las quejas procedentes, apoyando una atención equitativa y eficiente ante la creciente demanda de quejas que reciben diariamente.

\subsection{Perspectiva técnica}
Desde la perspectiva técnica, el desarrollo de este sistema se fundamenta en la adopción de una arquitectura de microservicios permitiendo la oportunidad de escalabilidad y mantenimiento independiente de sus componentes. Ya que la integración de tecnologías como Java Spring Boot y PostgreSQL permite el manejo seguro de información sensible, mientras que la incorporación de un clasificador basado en PLN demuestra una innovación clave para resolver los puntos significativos desde el área de primer contacto. Esta solución resuelve la fragmentación de datos y fortalece la integridad de un repositorio histórico de quejas para futuras consultas y auditorias.

\section{Objetivos}

\subsection{Objetivo General}
Desarrollar una plataforma web basada en una arquitectura de microservicios mediante la integración de módulos de gestión documental y un motor de procesamiento de lenguaje natural para sistematizar el flujo operativo de recepción, clasificación y seguimiento de quejas de la DDP.

\subsection{Objetivos específicos}
\begin{itemize}
    \item Modelar las etapas de proceso y resolución de quejas con base en las reglas de negocio detalladas en el Manual de Procedimientos de la DDP desde recepción inicial, clasificación, asignación, investigación, audiencia hasta resolución y archivo, para automatizar la transición de estados del expediente ("Pendiente", "Revisión", "Resuelto"), asegurando el cumplimiento estricto de términos legales como plazos de 2-5 días para primer contacto o 15 días para dictámenes preliminares bajo lo establecido en los Manuales de Prodecimientos de la Defensoria.
    \item Diseñar una arquitectura de microservicios desacoplada empleando el ecosistema de Java Spring Boot (incluyendo Spring Boot Actuator para endpoints y métricas, y Spring Security para autenticación) y comunicación estandarizada vía API REST (con soporte para HTTP/2, rate limiting y versioning semántico), para garantizar la independencia operativa total entre el sistema de gestión administrativa responsable de CRUD de expedientes, roles de usuario y flujos de aprobación y el módulo de inteligencia artificial encargado de PLN, embeddings y reglas simbólicas, permitiendo despliegues paralelos, escalado selectivo por demanda y resiliencia ante fallos aislados mediante patrones como Circuit Breaker (con Resilience4j). La implementación usará Docker para contenedorización, Kubernetes o Docker Swarm para orquestación con auto-scaling basado en CPU/memoria, colas de mensajes asíncronas (Apache Kafka para eventos de alto volumen como clasificaciones masivas), y un API Gateway central (Kong o Spring Cloud Gateway) para routing, CORS y throttling; se incorporará seguridad end-to-end con JWT/OAuth 2.0, encriptación TLS 1.3, observabilidad vía Jaeger para tracing distribuido, optimizando latencia en <200ms y soportando >1,000 RPS mientras cumple con regulaciones como la LFPDPPP.
    \item Implementar un clasificador de texto neuro-simbólico que combine Word Embeddings con sistemas basados en reglas expertas definidas por abogados de la DDP para determinar automáticamente la competencia de las solicitudes desde el primer contacto, categorizándolas por materias como derechos humanos, amparo, penal, civil o laboral, y así reducir la carga de clasificación manual en más del 60\%, liberando tiempo para análisis profundos. Este híbrido aprovecha la capacidad de los embeddings neuronales para capturar similitudes semánticas en lenguaje natural (incluso con variaciones dialectales o errores tipográficos comunes en quejas ciudadanas) y las reglas simbólicas para aplicar lógica deductiva precisa basada en normativas vigentes, con un flujo de fallback a supervisión humana para casos ambiguos (precisión inicial estimada en 85-92\%). El módulo se integrará vía API RESTful, permitirá entrenamiento continuo con datos anonimizados de expedientes reales y generará explicaciones interpretables (XAI) para justificar clasificaciones, asegurando transparencia y cumplimiento ético en la toma de decisiones asistida por IA.
    \item Construir un repositorio digital de evidencias y expedientes mediante el sistema de gestión de bases de datos PostgreSQL, optimizado para entornos de alta concurrencia y escalabilidad, con el fin de centralizar toda la documentación soporte incluyendo escaneos de documentos físicos, archivos multimedia, testimonios y peritajes y asegurar la integridad de la información jurídica manejada por la Defensoría mediante mecanismos como control de versiones, auditorías automáticas de cambios y respaldo en la nube con replicación en tiempo real. Este repositorio incorporará encriptación de datos en reposo y en tránsito (AES-256), control de acceso basado en roles (RBAC) alineado con las políticas internas de la DDP, y cumplimiento estricto de la LFPDPPP y normas ISO 27001 para protección de datos sensibles, facilitando búsquedas full-text con índices GIN, generación de hashes criptográficos para verificación de integridad y exportación segura de expedientes para juicios o auditorías externas, lo que minimiza pérdidas de información y agiliza el acceso en hasta un 70% comparado con sistemas manuales.
    \item Desarrollar un módulo de consulta y seguimiento para el quejoso a través de un sistema de folios digitales únicos (generados al momento del registro inicial) y notificaciones automáticas de estatus por correo electrónico, SMS o app móvil, brindando certidumbre en tiempo real sobre el avance del proceso desde la recepción hasta la resolución y reduciendo la saturación de los canales de comunicación presenciales en al menos un 50\%, según métricas de implementación similar en otras dependencias públicas. Este módulo incluirá un portal web accesible 24/7 con interfaz responsive, autenticación segura vía OTP (One-Time Password) y visualizaciones gráficas del timeline del expediente, permitiendo al usuario subir documentos adicionales, programar citas virtuales y recibir alertas personalizadas sobre plazos críticos, lo que optimiza la experiencia ciudadana y libera recursos del personal para tareas de alto valor.
\end{itemize}

\section{Alcances y Limitaciones}

\subsection{Alcances}
La plataforma web permitirá el registro, seguimiento y gestión integral de las quejas recibidas por la DDP, abarcando desde el primer contacto del ciudadano ya sea por vía digital, telefónica o presencial hasta la resolución final del expediente, incluyendo notificaciones automáticas y generación de reportes estadísticos en tiempo real. El sistema implementará un clasificador basado en Procesamiento de Lenguaje Natural (PLN) para asistir en la determinación automática de competencia de las solicitudes, categorizándolas por materias como derechos humanos, amparo, penal o civil, con una interfaz intuitiva que sugiera asignaciones preliminares y reduzca tiempos de procesamiento en hasta un 40\% según pruebas piloto. Asimismo, contará con un repositorio digital seguro para el almacenamiento de documentos y evidencias, cumpliendo estándares de encriptación AES-256 y autenticación multifactor, facilitando el acceso compartido y la trazabilidad de versiones para auditorías internas. Se desarrollarán módulos específicos y personalizados para cada rol institucional —Abogado Asesor, Subdefensor, Recepcionista, Área de Primer Contacto y Titular—, con dashboards adaptados que incluyan alertas prioritarias, flujos de trabajo colaborativos y herramientas de firma electrónica avanzada. Finalmente, se proporcionará un sistema de consulta pública mediante folios digitales únicos y QR codes, permitiendo a los usuarios externos rastrear el estatus de sus quejas de forma anónima y transparente, fomentando la confianza ciudadana en la institución.
\subsection{Limitaciones}
El sistema no sustituirá el criterio jurídico de los abogados de la DDP, sino que fungirá como una herramienta de apoyo en la clasificación inicial de quejas, permitiendo una triaje eficiente que acelere el flujo de trabajo sin comprometer la autonomía profesional. La plataforma no realizará resoluciones automáticas de quejas ni emitirá dictámenes legales vinculantes, ya que su diseño prioriza la revisión humana en todas las etapas críticas para garantizar el cumplimiento de normativas éticas y legales, como la protección de datos personales bajo la Ley Federal de Protección de Datos Personales en Posesión de los Particulares (LFPDPPP). Asimismo, el clasificador basado en IA requerirá un período de entrenamiento y ajuste continuo con datos reales anonimizados de casos históricos, lo que podría tomar varios meses para alcanzar niveles óptimos de precisión (por encima del 90\% en pruebas preliminares). En etapas iniciales, operará de manera asistida con supervisión humana obligatoria, donde los abogados validarán las clasificaciones sugeridas para mitigar sesgos algorítmicos, falsos positivos o negativos, y adaptarse progresivamente a la variabilidad del lenguaje jurídico en contextos mexicanos. Adicionalmente, la efectividad del sistema dependerá de la calidad y volumen de los datos de entrenamiento, por lo que se recomienda un protocolo de retroalimentación iterativa para su mejora continua.